\documentclass[11pt]{article}

\usepackage[utf8]{inputenc}
\usepackage[T1]{fontenc}
\usepackage{lmodern}
\usepackage{geometry}
\usepackage{amsmath,amssymb,amsfonts,amsthm,mathtools}
\usepackage{bm}
\usepackage{enumitem}
\usepackage{microtype}
\usepackage{hyperref}
\usepackage{titlesec}
\usepackage{booktabs}
\usepackage{array}
\usepackage{setspace}
\usepackage{mathrsfs}
\usepackage{physics}

\geometry{margin=1in}
\hypersetup{colorlinks=true, linkcolor=black, urlcolor=blue, citecolor=black}
\setlist[enumerate]{topsep=0.4em,itemsep=0.35em}
\setlist[itemize]{topsep=0.3em,itemsep=0.25em}
\titleformat{\section}{\large\bfseries}{\thesection.}{0.5em}{}
\titleformat{\subsection}{\normalsize\bfseries}{\thesubsection.}{0.5em}{}
\setstretch{1.06}

\newcommand{\Aether}{\AE{}ther}
\newcommand{\Aflow}{\AE{}ther-flow}
\newcommand{\TheoryName}{\Aflow{} Interpretation of Relativity}
\newcommand{\TheoryProgram}{\Aether{} / \Aflow{} framework}
\newcommand{\RespShapeReadout}{\widetilde q_{\mathrm{br}}}
\newcommand{\RespShapeCov}{\widetilde\kappa_{\mathrm{br}}}
\newcommand{\RespShapeRelax}{\widetilde\rho_{\mathrm{br}}}
\newcommand{\RespShapeWell}{\widetilde\lambda_{\mathrm{br}}}
\newcommand{\RespShapeEq}{\widetilde U_{\mathrm{br}}^{\ast}}
\newcommand{\RespShapeIso}{\widetilde\eta_{\mathrm{br}}}
\newcommand{\RespRawReadout}{\widehat q_{\mathrm{br}}}
\newcommand{\RespRawRef}{\widehat U_{\mathrm{br}}^{\mathrm{ref}}}
\newcommand{\RespRawCov}{\widehat\kappa_{\mathrm{br}}}
\newcommand{\RespRawRelax}{\widehat\rho_{\mathrm{br}}}
\newcommand{\RespRawWell}{\widehat\lambda_{\mathrm{br}}}
\newcommand{\RespRawEq}{\widehat U_{\mathrm{br}}^{\ast}}
\newcommand{\RespVacPhase}{\phi_{\mathrm{br}}}
\newcommand{\CurvProfileCan}{F_{\mathrm{br}}^{\mathrm{can}}}
\newcommand{\DownPkg}{\mathscr P_{\mathrm{down}}}
\newcommand{\ShapeTuple}{\mathfrak S_{\mathrm{br}}}
\newcommand{\ScaleMod}{s_{\mathrm{br}}}
\newcommand{\PrimClass}{\mathfrak C_{\mathrm{br}}}
\newcommand{\PrimRead}{r_{\mathrm{br}}}
\newcommand{\PrimCovSlot}{a_{\mathrm{br}}^{\varphi}}
\newcommand{\PrimRelaxSlot}{a_{\mathrm{br}}^{V}}
\newcommand{\PrimKinetic}{\bm a_{\mathrm{br}}}
\newcommand{\PrimWellSlot}{b_{\mathrm{br}}^{\lambda}}
\newcommand{\PrimEqSlot}{b_{\mathrm{br}}^{U}}
\newcommand{\PrimAnchor}{\bm b_{\mathrm{br}}}
\newcommand{\LiftSector}{\mathfrak D_{\mathrm{br}}^{\mathrm{amp}}}
\newcommand{\DeepReadAmp}{\xi_{\mathrm{br}}^{q}}
\newcommand{\DeepCovAmp}{\xi_{\mathrm{br}}^{\varphi}}
\newcommand{\DeepRelaxAmp}{\xi_{\mathrm{br}}^{V}}
\newcommand{\DeepKinAmp}{\bm\xi_{\mathrm{br}}}
\newcommand{\DeepWellAmp}{\zeta_{\mathrm{br}}^{\lambda}}
\newcommand{\DeepEqAmp}{\zeta_{\mathrm{br}}^{U}}
\newcommand{\DeepAnchorAmp}{\bm\zeta_{\mathrm{br}}}
\newcommand{\DeepScaleAmp}{\varsigma_{\mathrm{br}}}
\newcommand{\AmpMap}{\Pi_{\mathrm{amp}}}
\newcommand{\PrimMap}{\Pi_{\mathrm{shape}}}
\newcommand{\OriginSector}{\mathfrak D_{\mathrm{br}}^{\mathrm{ori}}}
\newcommand{\AbsAmpSector}{\mathfrak M_{\mathrm{br}}}
\newcommand{\PolaritySector}{\mathfrak P_{\mathrm{br}}}
\newcommand{\AbsReadAmp}{\chi_{\mathrm{br}}^{q}}
\newcommand{\AbsCovAmp}{\chi_{\mathrm{br}}^{\varphi}}
\newcommand{\AbsRelaxAmp}{\chi_{\mathrm{br}}^{V}}
\newcommand{\AbsKinAmp}{\bm\chi_{\mathrm{br}}}
\newcommand{\AbsWellAmp}{\upsilon_{\mathrm{br}}^{\lambda}}
\newcommand{\AbsEqAmp}{\upsilon_{\mathrm{br}}^{U}}
\newcommand{\AbsAnchorAmp}{\bm\upsilon_{\mathrm{br}}}
\newcommand{\AbsScaleAmp}{\omega_{\mathrm{br}}}
\newcommand{\PolRead}{\epsilon_{\mathrm{br}}^{q}}
\newcommand{\PolCov}{\epsilon_{\mathrm{br}}^{\varphi}}
\newcommand{\PolRelax}{\epsilon_{\mathrm{br}}^{V}}
\newcommand{\PolWell}{\delta_{\mathrm{br}}^{\lambda}}
\newcommand{\PolEq}{\delta_{\mathrm{br}}^{U}}
\newcommand{\PolScale}{\sigma_{\mathrm{br}}}
\newcommand{\OriginMap}{\Pi_{\mathrm{ori}}}
\newcommand{\DeepMap}{\Pi_{\mathrm{deep}}}
\newcommand{\PairSector}{\mathfrak D_{\mathrm{br}}^{\mathrm{pair}}}
\newcommand{\PairMap}{\Pi_{\mathrm{pair}}}
\newcommand{\PairFiber}{\sim_{\mathrm{fib}}}
\newcommand{\PairReadPlus}{Q_{\mathrm{br}}^{+}}
\newcommand{\PairReadMinus}{Q_{\mathrm{br}}^{-}}
\newcommand{\PairCovPlus}{\Phi_{\mathrm{br}}^{+}}
\newcommand{\PairCovMinus}{\Phi_{\mathrm{br}}^{-}}
\newcommand{\PairRelaxPlus}{\mathcal V_{\mathrm{br}}^{+}}
\newcommand{\PairRelaxMinus}{\mathcal V_{\mathrm{br}}^{-}}
\newcommand{\PairWellPlus}{\Lambda_{\mathrm{br}}^{+}}
\newcommand{\PairWellMinus}{\Lambda_{\mathrm{br}}^{-}}
\newcommand{\PairEqPlus}{\mathcal U_{\mathrm{br}}^{+}}
\newcommand{\PairEqMinus}{\mathcal U_{\mathrm{br}}^{-}}
\newcommand{\PairScalePlus}{\Omega_{\mathrm{br}}^{+}}
\newcommand{\PairScaleMinus}{\Omega_{\mathrm{br}}^{-}}
\newcommand{\GapRead}{\Delta Q_{\mathrm{br}}}
\newcommand{\GapCov}{\Delta \Phi_{\mathrm{br}}}
\newcommand{\GapRelax}{\Delta \mathcal V_{\mathrm{br}}}
\newcommand{\GapWell}{\Delta \Lambda_{\mathrm{br}}}
\newcommand{\GapEq}{\Delta \mathcal U_{\mathrm{br}}}
\newcommand{\GapScale}{\Delta \Omega_{\mathrm{br}}}
\newcommand{\Rstar}{\mathbb R^{\times}}

\newtheorem{definition}{Definition}
\newtheorem{proposition}{Proposition}
\newtheorem{remark}{Remark}
\newtheorem{corollary}{Corollary}
\newtheorem{theorem}{Theorem}

\title{\textbf{The \TheoryName{}}\\[0.4em]
\large Nonlocal Bridge Self-Response Off-Shell Profile-Selection Bridge-Order Orbit-Shape/Modulus Primitive-Class Absolute-Amplitude/Polarity Origin Note}
\author{}
\date{}

\begin{document}

\maketitle

\begin{abstract}
The benchmark exact-theory statement of \TheoryName{} remains fixed and is not reopened here. The overview-first exact-closure package is still the public front door. The present note acts only on the optional deeper derivational bridge, and it begins from the now-recorded amplitude-sector origin note immediately below the primitive-class amplitude-lift layer.

That amplitude-sector note already pushed the derivational line one step deeper in a clean way. It replaced the six signed amplitudes by the oriented absolute-amplitude sector
\[
\OriginSector
=
\AbsAmpSector\times\PolaritySector,
\]
with positive absolute amplitudes
\[
(\AbsReadAmp,\AbsCovAmp,\AbsRelaxAmp,\AbsWellAmp,\AbsEqAmp,\AbsScaleAmp)
\]
and discrete polarity data
\[
(\PolRead,\PolCov,\PolRelax,\PolWell,\PolEq,\PolScale),
\]
and it showed that the same signed amplitudes, the same primitive class, the same raw-orbit lift, and the same downstream package were recovered exactly. One narrower burden remained. The positive absolute amplitudes and their polarity data were still left as the deepest disciplined layer.

This note chooses a concrete still-deeper realization of that layer. The new deeper sector is the unequal positive-pair sector
\[
\PairSector\subset(\mathbb R_{>0}^{2})^{6},
\]
with coordinates
\[
(\PairReadPlus,\PairReadMinus,\PairCovPlus,\PairCovMinus,\PairRelaxPlus,\PairRelaxMinus,
\PairWellPlus,\PairWellMinus,\PairEqPlus,\PairEqMinus,\PairScalePlus,\PairScaleMinus),
\]
subject only to nonzero pair gaps. The map into \(\OriginSector\) is componentwise gap magnitude and ordering sign:
\[
\AbsReadAmp=\abs{\PairReadPlus-\PairReadMinus},
\qquad
\PolRead=\operatorname{sgn}(\PairReadPlus-\PairReadMinus),
\]
and similarly in the other five slots. Thus each positive absolute amplitude is the size of an unequal positive-pair gap, while each discrete polarity is the order orientation of that same gap.

Using the frozen maps exactly as already recorded, the note then proves four facts. First, every oriented absolute-amplitude datum admits many lifts to \(\PairSector\), explicitly parameterized by positive midpoint data, so the absolute-amplitude/polarity layer is no longer primitive either. Second, the composite map into the signed amplitudes is simply the six pair gaps themselves, and therefore the frozen amplitude lift reproduces the same primitive class, the same readout block, the same ordered kinetic pair, the same radial anchor pair, and the same positive orbit modulus. Third, the same raw-orbit lift, the same downstream density/current block, the same primitive bridge-order block, the same phase-neutral minimizing branch, the same canonical profile, and the same dressed bridge map are recovered unchanged. Fourth, the only new quotient visible at this level is the internal pair-fiber quotient identifying different unequal positive-pair realizations of the same oriented absolute-amplitude datum. That quotient does not identify distinct positive moduli and it introduces no vacuum-angle locking or angle quotient. Therefore the disciplined current verdict remains unchanged: the positive scale orbit is still a canonical-normalization orbit rather than a proved redundancy, and the global \(O(2)\) vacuum angle remains residual.

The scope is deliberately narrow. This note does not claim that the unequal positive-pair sector is the unique final substrate microphysics. Its narrower positive result is that the positive absolute amplitudes and discrete polarity data feeding the amplitude-sector note are now themselves generated from a still deeper explicit sector, again without disturbing the already-positive downstream package.
\end{abstract}

\section{Introduction}

The active repository no longer needs another bridge note in order to justify its public theory claim. The exact-closure sequence already presents \TheoryName{} as the disciplined exact relativistic theory statement built on the \Aether{} / \Aflow{} ontology, and that benchmark package remains the front-door reading order. The present note therefore does not strengthen the flagship theory claim. It acts only on the optional deeper derivational continuation beneath that fixed benchmark.

On that optional continuation line the latest burden is now narrower again. The primitive-class amplitude-lift note stopped treating the positive primitive coefficients as terminal input by realizing them as squares of six signed amplitudes. The amplitude-sector origin note then stopped treating those signed amplitudes as terminal input by realizing them as positive absolute amplitudes together with discrete polarity data. What remains open after that note is still sharper: whether those positive absolute amplitudes and polarity data can themselves be generated from still deeper substrate variables without disturbing the frozen downstream chain.

That is the exact burden of the present manuscript. It does not reopen another same-layer discussion about signed amplitudes, it does not alter the already-positive raw-orbit lift, and it does not reopen the current scale-orbit or residual-angle verdicts unless the new deeper sector genuinely forces that revision. Instead it chooses the smallest concrete deeper realization in which positive absolute amplitudes arise as gap magnitudes of unequal positive pairs and the polarity data arise as their order orientation, and it then checks that the entire already-frozen path into the signed amplitudes, the primitive class, the orbit-shape/modulus data, the raw orbit, and the downstream package remains unchanged.

\paragraph{Framework and claim boundary.}
\TheoryProgram{} denotes the broader \Aether{} / \Aflow{} framework built on the claims that the \Aether{} is the underlying four-dimensional substrate of reality, the \Aflow{} is its intrinsic ordered motion, observed three-dimensional space is the local experiential slice of that deeper substrate, S-time is the experienced order of change arising from matter, light, and the \Aflow{}, and observed expansion is the three-dimensional appearance of deeper four-dimensional ordered motion. Within that framework, \TheoryName{} denotes the exact-closure theory in which the effective gravitational dynamics are adopted to be exactly Einsteinian with universal matter coupling. In the active sequence, \emph{adoption} means use of that established relativistic dynamics without claiming substrate derivation, while \emph{derivation} is reserved for a first-principles recovery from explicit substrate variables.

This note stays strictly inside that boundary. Its positive claim is not that the unique final microscopic completion has already been isolated. Its positive claim is narrower: the positive absolute-amplitude block and the discrete polarity data feeding the amplitude-sector note are no longer terminal input either. They arise from an explicit still-deeper unequal positive-pair sector through a frozen map that preserves the already-positive downstream package exactly.

\section{Frozen Amplitude-Sector Layer}

\begin{definition}[Frozen oriented absolute-amplitude sector]
\label{def:frozenorigin}
The amplitude-sector origin note fixed the oriented absolute-amplitude sector
\begin{equation}
\OriginSector
:=
\AbsAmpSector\times\PolaritySector,
\label{eq:originsector}
\end{equation}
where
\begin{equation}
\AbsAmpSector
:=
\mathbb R_{>0}\times(\mathbb R_{>0})^{2}\times(\mathbb R_{>0})^{2}\times\mathbb R_{>0}
\label{eq:abssector}
\end{equation}
has coordinates
\begin{equation}
(\AbsReadAmp,\AbsKinAmp,\AbsAnchorAmp,\AbsScaleAmp)
=
(\AbsReadAmp,(\AbsCovAmp,\AbsRelaxAmp),(\AbsWellAmp,\AbsEqAmp),\AbsScaleAmp),
\label{eq:abscoords}
\end{equation}
and
\begin{equation}
\PolaritySector
:=
\{\pm1\}\times\{\pm1\}^{2}\times\{\pm1\}^{2}\times\{\pm1\}
\label{eq:polsector}
\end{equation}
has coordinates
\begin{equation}
(\PolRead,\PolCov,\PolRelax,\PolWell,\PolEq,\PolScale).
\label{eq:polcoords}
\end{equation}
The frozen map into the signed amplitude sector
\begin{equation}
\LiftSector
:=
\Rstar\times(\Rstar)^{2}\times(\Rstar)^{2}\times\Rstar
\label{eq:liftsector}
\end{equation}
is
\begin{equation}
\OriginMap:\OriginSector\longrightarrow\LiftSector
\label{eq:originmap}
\end{equation}
with
\begin{align}
\OriginMap(\AbsReadAmp,\AbsKinAmp,\AbsAnchorAmp,\AbsScaleAmp;\PolRead,\PolCov,\PolRelax,\PolWell,\PolEq,\PolScale)
=\;&
(\PolRead\AbsReadAmp,\PolCov\AbsCovAmp,\PolRelax\AbsRelaxAmp,
\nonumber\\
&\PolWell\AbsWellAmp,\PolEq\AbsEqAmp,\PolScale\AbsScaleAmp).
\label{eq:originmapexplicit}
\end{align}
\end{definition}

\begin{definition}[Frozen maps below the amplitude-sector layer]
\label{def:frozenmaps}
The frozen amplitude lift into the primitive class
\begin{equation}
\PrimClass
:=
(\PrimRead,\PrimKinetic,\PrimAnchor,\ScaleMod)
\in
\mathbb R_{>0}\times\mathbb R_{>0}^{2}\times\mathbb R_{>0}^{2}\times\mathbb R_{>0}
\label{eq:primclass}
\end{equation}
is
\begin{equation}
\AmpMap:\LiftSector\longrightarrow\PrimClass
\label{eq:ampmap}
\end{equation}
given by
\begin{equation}
\AmpMap(\DeepReadAmp,\DeepKinAmp,\DeepAnchorAmp,\DeepScaleAmp)
=
\bigl(
(\DeepReadAmp)^{2},
((\DeepCovAmp)^{2},(\DeepRelaxAmp)^{2}),
((\DeepWellAmp)^{2},(\DeepEqAmp)^{2}),
(\DeepScaleAmp)^{2}
\bigr).
\label{eq:ampmapexplicit}
\end{equation}
The frozen primitive-class setup map is
\begin{equation}
\PrimMap(\PrimRead,\PrimKinetic,\PrimAnchor,\ScaleMod)
=
(\RespShapeReadout,\RespShapeCov,\RespShapeRelax,\RespShapeWell,\RespShapeEq,\ScaleMod)
\label{eq:primmap}
\end{equation}
with
\begin{equation}
\RespShapeReadout=\PrimRead,
\qquad
\RespShapeCov=\PrimCovSlot,
\qquad
\RespShapeRelax=\PrimRelaxSlot,
\label{eq:shapeA}
\end{equation}
\begin{equation}
\RespShapeWell=\PrimWellSlot,
\qquad
\RespShapeEq=\PrimEqSlot.
\label{eq:shapeB}
\end{equation}
The frozen raw-orbit lift remains
\begin{equation}
\RespRawReadout=\ScaleMod\,\RespShapeReadout,
\qquad
\RespRawRef=\ScaleMod,
\qquad
\RespRawCov=\ScaleMod^{2}\RespShapeCov,
\label{eq:rawliftA}
\end{equation}
\begin{equation}
\RespRawRelax=\ScaleMod^{2}\RespShapeRelax,
\qquad
\RespRawWell=\RespShapeWell,
\qquad
\RespRawEq=\ScaleMod\,\RespShapeEq.
\label{eq:rawliftB}
\end{equation}
\end{definition}

\begin{definition}[Frozen downstream package and verdict criterion]
\label{def:frozendown}
At the present derivational stage, the downstream package \(\DownPkg\) means the same already-recorded density/current block, the same primitive bridge-order block, the same phase-neutral minimizing branch, the same canonical profile
\begin{equation}
\CurvProfileCan(x)
=
1+\RespShapeEq\,x^{2}
=
1+\frac{\RespRawEq}{\RespRawRef}x^{2},
\label{eq:profile}
\end{equation}
and the same dressed bridge map, all understood to depend on deeper data only through the induced pair \((\ShapeTuple,\ScaleMod)\), the raw lift \eqref{eq:rawliftA}--\eqref{eq:rawliftB}, and the same minimizing branch.

At the same disciplined scope, a deeper continuation changes the current scale-orbit or residual vacuum-angle verdicts only if it provides at least one of the following:
\begin{enumerate}
    \item a quotient identifying distinct positive values of \(\AbsScaleAmp\), or equivalently distinct positive values of \(\ScaleMod=(\AbsScaleAmp)^{2}\);
    \item a locking mechanism or true quotient identifying distinct constant vacuum angles \(\RespVacPhase=\phi_{0}\).
\end{enumerate}
\end{definition}

\begin{remark}
Definitions \ref{def:frozenorigin}--\ref{def:frozendown} are frozen input in this note. The present manuscript does not modify the amplitude-sector realization or any already-positive downstream formula. It acts only below that layer.
\end{remark}

\section{A Still-Deeper Unequal Positive-Pair Sector}

\begin{definition}[Unequal positive-pair sector]
\label{def:pairsector}
Write points of \((\mathbb R_{>0}^{2})^{6}\) as
\begin{equation}
z=
(\PairReadPlus,\PairReadMinus,\PairCovPlus,\PairCovMinus,\PairRelaxPlus,\PairRelaxMinus,
\PairWellPlus,\PairWellMinus,\PairEqPlus,\PairEqMinus,\PairScalePlus,\PairScaleMinus).
\label{eq:paircoords}
\end{equation}
Define the deeper sector
\begin{equation}
\PairSector
:=
\left\{
z\in(\mathbb R_{>0}^{2})^{6}
\;:\;
\GapRead\neq0,\ \GapCov\neq0,\ \GapRelax\neq0,\ \GapWell\neq0,\ \GapEq\neq0,\ \GapScale\neq0
\right\},
\label{eq:pairsector}
\end{equation}
where the six pair gaps are
\begin{equation}
\begin{aligned}
\GapRead&:=\PairReadPlus-\PairReadMinus,\\
\GapCov&:=\PairCovPlus-\PairCovMinus,\\
\GapRelax&:=\PairRelaxPlus-\PairRelaxMinus,
\end{aligned}
\label{eq:gapsA}
\end{equation}
\begin{equation}
\begin{aligned}
\GapWell&:=\PairWellPlus-\PairWellMinus,\\
\GapEq&:=\PairEqPlus-\PairEqMinus,\\
\GapScale&:=\PairScalePlus-\PairScaleMinus.
\end{aligned}
\label{eq:gapsB}
\end{equation}
\end{definition}

\begin{definition}[Gap/order map into the oriented absolute-amplitude sector]
\label{def:pairmap}
Define
\begin{equation}
\PairMap:\PairSector\longrightarrow\OriginSector
\label{eq:pairmap}
\end{equation}
componentwise by
\begin{equation}
\AbsReadAmp=\abs{\GapRead},
\qquad
\AbsCovAmp=\abs{\GapCov},
\qquad
\AbsRelaxAmp=\abs{\GapRelax},
\label{eq:pairabsA}
\end{equation}
\begin{equation}
\AbsWellAmp=\abs{\GapWell},
\qquad
\AbsEqAmp=\abs{\GapEq},
\qquad
\AbsScaleAmp=\abs{\GapScale},
\label{eq:pairabsB}
\end{equation}
and
\begin{equation}
\PolRead=\operatorname{sgn}(\GapRead),
\qquad
\PolCov=\operatorname{sgn}(\GapCov),
\qquad
\PolRelax=\operatorname{sgn}(\GapRelax),
\label{eq:pairpolA}
\end{equation}
\begin{equation}
\PolWell=\operatorname{sgn}(\GapWell),
\qquad
\PolEq=\operatorname{sgn}(\GapEq),
\qquad
\PolScale=\operatorname{sgn}(\GapScale).
\label{eq:pairpolB}
\end{equation}
\end{definition}

\begin{proposition}[The pair map lands in the frozen oriented sector]
\label{prop:lands}
The map \(\PairMap\) is well defined. Every image point consists of a positive absolute-amplitude block together with a discrete polarity block.
\end{proposition}

\begin{proof}
By Definition \ref{def:pairsector}, each gap in \eqref{eq:gapsA}--\eqref{eq:gapsB} is nonzero. Hence each absolute value in \eqref{eq:pairabsA}--\eqref{eq:pairabsB} is strictly positive, and each sign in \eqref{eq:pairpolA}--\eqref{eq:pairpolB} is an element of \(\{\pm1\}\). Therefore \(\PairMap\) indeed takes values in \(\AbsAmpSector\times\PolaritySector=\OriginSector\).
\end{proof}

\begin{proposition}[Every oriented absolute-amplitude datum admits many positive-pair lifts]
\label{prop:surjective}
The map \(\PairMap\) is surjective. More precisely, let
\[
(\AbsReadAmp,\AbsCovAmp,\AbsRelaxAmp,\AbsWellAmp,\AbsEqAmp,\AbsScaleAmp;
\PolRead,\PolCov,\PolRelax,\PolWell,\PolEq,\PolScale)\in\OriginSector.
\]
Choose any midpoint tuple
\[
(C_{q},C_{\varphi},C_{V},C_{\lambda},C_{U},C_{\omega})
\in(\mathbb R_{>0})^{6}
\]
satisfying
\[
C_{q}>\frac{\AbsReadAmp}{2},
\quad
C_{\varphi}>\frac{\AbsCovAmp}{2},
\quad
C_{V}>\frac{\AbsRelaxAmp}{2},
\quad
C_{\lambda}>\frac{\AbsWellAmp}{2},
\quad
C_{U}>\frac{\AbsEqAmp}{2},
\quad
C_{\omega}>\frac{\AbsScaleAmp}{2}.
\]
Then the formulas
\begin{align}
\PairReadPlus&=C_{q}+\frac{1}{2}\PolRead\AbsReadAmp,
&
\PairReadMinus&=C_{q}-\frac{1}{2}\PolRead\AbsReadAmp,
\nonumber\\
\PairCovPlus&=C_{\varphi}+\frac{1}{2}\PolCov\AbsCovAmp,
&
\PairCovMinus&=C_{\varphi}-\frac{1}{2}\PolCov\AbsCovAmp,
\nonumber\\
\PairRelaxPlus&=C_{V}+\frac{1}{2}\PolRelax\AbsRelaxAmp,
&
\PairRelaxMinus&=C_{V}-\frac{1}{2}\PolRelax\AbsRelaxAmp,
\label{eq:liftsA}
\end{align}
\begin{align}
\PairWellPlus&=C_{\lambda}+\frac{1}{2}\PolWell\AbsWellAmp,
&
\PairWellMinus&=C_{\lambda}-\frac{1}{2}\PolWell\AbsWellAmp,
\nonumber\\
\PairEqPlus&=C_{U}+\frac{1}{2}\PolEq\AbsEqAmp,
&
\PairEqMinus&=C_{U}-\frac{1}{2}\PolEq\AbsEqAmp,
\nonumber\\
\PairScalePlus&=C_{\omega}+\frac{1}{2}\PolScale\AbsScaleAmp,
&
\PairScaleMinus&=C_{\omega}-\frac{1}{2}\PolScale\AbsScaleAmp
\label{eq:liftsB}
\end{align}
define a point of \(\PairSector\) mapping back to the chosen element of \(\OriginSector\).
\end{proposition}

\begin{proof}
Because each midpoint exceeds half the corresponding positive absolute amplitude, every entry defined by \eqref{eq:liftsA}--\eqref{eq:liftsB} is strictly positive. Moreover,
\[
\PairReadPlus-\PairReadMinus=\PolRead\AbsReadAmp,
\qquad
\PairCovPlus-\PairCovMinus=\PolCov\AbsCovAmp,
\]
and similarly in the remaining slots. Therefore the gaps are all nonzero, so the resulting tuple lies in \(\PairSector\). Applying \(\PairMap\) then recovers exactly the original absolute amplitudes and polarity labels. Since the midpoint tuple is arbitrary subject to the positivity inequalities, the lift is not unique.
\end{proof}

\begin{corollary}[The absolute-amplitude/polarity layer is not primitive]
\label{cor:notprimitive}
At the level of the present note, the oriented absolute-amplitude sector \(\OriginSector\) is not primitive. It is the image of the still-deeper unequal positive-pair sector \(\PairSector\) under the map \(\PairMap\).
\end{corollary}

\begin{proof}
This is exactly Proposition \ref{prop:surjective}.
\end{proof}

\begin{remark}
The extra data introduced by Proposition \ref{prop:surjective} are the midpoint or offset variables \(C_{q},C_{\varphi},C_{V},C_{\lambda},C_{U},C_{\omega}\). The present note does not yet derive or select them. It only shows that the positive absolute amplitudes and polarity data already arise from a concrete deeper layer with a transparent internal fiber.
\end{remark}

\section{Composite Map Into the Signed Amplitudes and Frozen Interface}

\begin{theorem}[The signed amplitudes are exactly the pair gaps]
\label{thm:signed}
The composite map \(\OriginMap\circ\PairMap:\PairSector\to\LiftSector\) is
\begin{equation}
(\OriginMap\circ\PairMap)
=
(\GapRead,\GapCov,\GapRelax,\GapWell,\GapEq,\GapScale).
\label{eq:signedgaps}
\end{equation}
Equivalently,
\begin{equation}
\DeepReadAmp=\GapRead,
\qquad
\DeepCovAmp=\GapCov,
\qquad
\DeepRelaxAmp=\GapRelax,
\label{eq:deepgapsA}
\end{equation}
\begin{equation}
\DeepWellAmp=\GapWell,
\qquad
\DeepEqAmp=\GapEq,
\qquad
\DeepScaleAmp=\GapScale.
\label{eq:deepgapsB}
\end{equation}
\end{theorem}

\begin{proof}
Each component of \(\OriginMap\circ\PairMap\) is of the form
\[
\operatorname{sgn}(\Delta)\abs{\Delta}=\Delta
\]
for the corresponding nonzero gap \(\Delta\). Substituting \eqref{eq:pairabsA}--\eqref{eq:pairpolB} into \eqref{eq:originmapexplicit} therefore yields exactly \eqref{eq:signedgaps}--\eqref{eq:deepgapsB}.
\end{proof}

\begin{definition}[Composite map into the primitive class and orbit-shape/modulus data]
\label{def:deepmap}
Define
\begin{equation}
\DeepMap
:=
\PrimMap\circ\AmpMap\circ\OriginMap\circ\PairMap.
\label{eq:deepmap}
\end{equation}
Explicitly, the primitive-class slots become
\begin{equation}
\PrimRead=(\GapRead)^{2},
\qquad
\PrimCovSlot=(\GapCov)^{2},
\qquad
\PrimRelaxSlot=(\GapRelax)^{2},
\label{eq:primgapsA}
\end{equation}
\begin{equation}
\PrimWellSlot=(\GapWell)^{2},
\qquad
\PrimEqSlot=(\GapEq)^{2},
\qquad
\ScaleMod=(\GapScale)^{2},
\label{eq:primgapsB}
\end{equation}
and hence
\begin{equation}
\RespShapeReadout=(\GapRead)^{2},
\qquad
\RespShapeCov=(\GapCov)^{2},
\qquad
\RespShapeRelax=(\GapRelax)^{2},
\label{eq:shapegapsA}
\end{equation}
\begin{equation}
\RespShapeWell=(\GapWell)^{2},
\qquad
\RespShapeEq=(\GapEq)^{2},
\qquad
\ScaleMod=(\GapScale)^{2}.
\label{eq:shapegapsB}
\end{equation}
\end{definition}

\begin{corollary}[Same readout block, ordered kinetic pair, radial anchor pair, and positive orbit modulus]
\label{cor:blockmap}
Under \(\DeepMap\), the frozen interface blocks are recovered as
\begin{enumerate}
    \item the readout block
    \[
    \RespShapeReadout=(\GapRead)^{2};
    \]
    \item the ordered kinetic pair
    \[
    (\RespShapeCov,\RespShapeRelax)=((\GapCov)^{2},(\GapRelax)^{2});
    \]
    \item the radial anchor pair
    \[
    (\RespShapeWell,\RespShapeEq)=((\GapWell)^{2},(\GapEq)^{2});
    \]
    \item the positive orbit modulus
    \[
    \ScaleMod=(\GapScale)^{2}.
    \]
\end{enumerate}
\end{corollary}

\begin{proof}
This is exactly the explicit content of Definition \ref{def:deepmap}.
\end{proof}

\begin{remark}
Theorem \ref{thm:signed} and Corollary \ref{cor:blockmap} isolate the exact gain of the present note. The signed amplitudes are no longer primitive, the positive absolute amplitudes are no longer primitive, and yet the frozen interface remains exactly the same because it sees only the six nonzero pair gaps through their squares.
\end{remark}

\section{Same Raw-Orbit Lift and Same Downstream Package}

\begin{theorem}[The raw-orbit lift is unchanged]
\label{thm:rawlift}
Composing \(\DeepMap\) with the frozen raw lift \eqref{eq:rawliftA}--\eqref{eq:rawliftB} yields exactly the same raw orbit as in the amplitude-lift and amplitude-sector notes. In the present variables,
\begin{equation}
\RespRawReadout=(\GapScale\,\GapRead)^{2},
\qquad
\RespRawRef=(\GapScale)^{2},
\qquad
\RespRawCov=(\GapScale\,\GapCov)^{2},
\label{eq:rawgapsA}
\end{equation}
\begin{equation}
\RespRawRelax=(\GapScale\,\GapRelax)^{2},
\qquad
\RespRawWell=(\GapWell)^{2},
\qquad
\RespRawEq=(\GapScale\,\GapEq)^{2}.
\label{eq:rawgapsB}
\end{equation}
\end{theorem}

\begin{proof}
Insert \eqref{eq:shapegapsA}--\eqref{eq:shapegapsB} into the raw lift \eqref{eq:rawliftA}--\eqref{eq:rawliftB}. This gives
\[
\RespRawReadout=(\GapScale)^{2}(\GapRead)^{2}=(\GapScale\,\GapRead)^{2},
\qquad
\RespRawRef=(\GapScale)^{2},
\]
\[
\RespRawCov=(\GapScale)^{2}(\GapCov)^{2}=(\GapScale\,\GapCov)^{2},
\qquad
\RespRawRelax=(\GapScale)^{2}(\GapRelax)^{2}=(\GapScale\,\GapRelax)^{2},
\]
\[
\RespRawWell=(\GapWell)^{2},
\qquad
\RespRawEq=(\GapScale)^{2}(\GapEq)^{2}=(\GapScale\,\GapEq)^{2}.
\]
These are exactly \eqref{eq:rawgapsA}--\eqref{eq:rawgapsB}. By Theorem \ref{thm:signed}, they are the same formulas already recorded one layer up, rewritten in terms of pair gaps.
\end{proof}

\begin{definition}[Induced orbit-level isotropic coefficient]
\label{def:isocoeff}
The induced orbit-level isotropic coefficient is
\begin{equation}
\RespShapeIso
:=
\frac{\RespShapeCov\,\RespShapeRelax}{\RespShapeCov+\RespShapeRelax}
=
\frac{(\GapCov)^{2}(\GapRelax)^{2}}{(\GapCov)^{2}+(\GapRelax)^{2}}.
\label{eq:iso}
\end{equation}
\end{definition}

\begin{theorem}[Same phase-neutral minimizing branch]
\label{thm:branch}
The orbit-level energy induced by Definition \ref{def:deepmap} is the same frozen-interface energy as in the setup, amplitude-lift, and amplitude-sector notes, now with coefficients \eqref{eq:shapegapsA}--\eqref{eq:shapegapsB}. Its distinguished minimizing constant branch is therefore
\begin{equation}
\ScaleMod(x)\equiv(\GapScale)^{2},
\qquad
U(x)\equiv(\GapScale\,\GapEq)^{2},
\qquad
V(x)\equiv0,
\qquad
\RespVacPhase(x)\equiv\phi_{0},
\label{eq:branch}
\end{equation}
for arbitrary constant \(\phi_{0}\in\mathbb R/2\pi\mathbb Z\). In particular, the minimizing branch remains phase-neutral.
\end{theorem}

\begin{proof}
Definition \ref{def:deepmap} produces exactly the same orbit-shape coefficients and positive modulus as the amplitude-sector note, because those coefficients depended there only on the signed amplitudes through their squares and Theorem \ref{thm:signed} identifies those signed amplitudes with the present pair gaps. Theorem \ref{thm:rawlift} shows that the raw-orbit lift is also unchanged. The minimizing-family computation in the frozen orbit-level energy therefore gives exactly the same constant branch as before, which becomes \eqref{eq:branch} after substituting the present gap formulas. Because \(\phi_{0}\) remains arbitrary and constant while \(V\equiv0\), the branch is still phase-neutral.
\end{proof}

\begin{theorem}[Same downstream density/current block, primitive bridge-order block, canonical profile, and dressed bridge map]
\label{thm:downstream}
For the deeper sector \(\PairSector\), the downstream package \(\DownPkg\) is recovered unchanged. In particular:
\begin{enumerate}
    \item the same downstream density/current block is recovered;
    \item the same primitive bridge-order block is recovered;
    \item the same phase-neutral minimizing branch of Theorem \ref{thm:branch} is recovered;
    \item the same canonical profile is recovered, now written as
    \begin{equation}
    \CurvProfileCan(x)
    =
    1+\RespShapeEq\,x^{2}
    =
    1+(\GapEq)^{2}x^{2}
    =
    1+\frac{\RespRawEq}{\RespRawRef}x^{2};
    \label{eq:profilegaps}
    \end{equation}
    \item the same dressed bridge map is recovered.
\end{enumerate}
\end{theorem}

\begin{proof}
Definition \ref{def:deepmap} shows that the present deeper sector feeds the frozen interface only through the induced pair \((\ShapeTuple,\ScaleMod)\). Theorem \ref{thm:rawlift} shows that the same raw tuple is then generated. Theorem \ref{thm:branch} shows that the same phase-neutral minimizing branch survives. By Definition \ref{def:frozendown}, all recorded downstream constructions depend on deeper data only through exactly those objects. Therefore every component of \(\DownPkg\) remains unchanged.

The canonical profile formula \eqref{eq:profilegaps} is immediate from \(\RespShapeEq=(\GapEq)^{2}\) and \(\RespRawEq/\RespRawRef=\RespShapeEq\).
\end{proof}

\begin{remark}
Theorem \ref{thm:downstream} is the exact reason the present note is honest. It deepens the origin of the absolute-amplitude/polarity layer itself without trying to earn novelty by rewriting formulas that are already fixed positively downstream.
\end{remark}

\section{What Changes and What Does Not}

\begin{definition}[Pair-fiber relation]
\label{def:fiber}
Define an equivalence relation \(\PairFiber\) on \(\PairSector\) by
\begin{equation}
z\PairFiber z'
\quad\Longleftrightarrow\quad
\PairMap(z)=\PairMap(z')
\label{eq:fiber}
\end{equation}
for \(z,z'\in\PairSector\).
\end{definition}

\begin{proposition}[The pair-fiber quotient is exactly the oriented absolute-amplitude sector]
\label{prop:fiberquot}
The quotient of \(\PairSector\) by the pair-fiber relation is canonically identified with \(\OriginSector\):
\begin{equation}
\PairSector/\PairFiber
\;\cong\;
\OriginSector.
\label{eq:fiberquot}
\end{equation}
\end{proposition}

\begin{proof}
By definition, two points of \(\PairSector\) are \(\PairFiber\)-equivalent exactly when they have the same image under \(\PairMap\). Therefore the quotient by \(\PairFiber\) is canonically the image of \(\PairMap\). Proposition \ref{prop:surjective} shows that the image is all of \(\OriginSector\). Hence \eqref{eq:fiberquot} holds.
\end{proof}

\begin{remark}
The pair-fiber quotient is an internal deeper redundancy only at the level of generating the already-recorded oriented absolute-amplitude data. It is not yet claimed to be a physically selected quotient or a uniquely dynamical gauge symmetry of the substrate theory.
\end{remark}

\begin{theorem}[The present deeper sector introduces no new verdict-changing structure]
\label{thm:verdict}
The unequal positive-pair sector \(\PairSector\) introduces no new verdict-changing structure in the sense of Definition \ref{def:frozendown}. Therefore:
\begin{enumerate}
    \item the positive scale orbit remains a canonical-normalization orbit rather than a proved redundancy;
    \item the global \(O(2)\) vacuum angle remains residual.
\end{enumerate}
\end{theorem}

\begin{proof}
By Proposition \ref{prop:fiberquot}, the only canonical quotient introduced at the present deeper level is the pair-fiber quotient collapsing distinct positive-pair realizations of the same oriented absolute-amplitude datum. That quotient does not identify distinct positive values of \(\AbsScaleAmp=\abs{\GapScale}\). Equivalently, it does not identify distinct positive values of the orbit modulus \(\ScaleMod=(\GapScale)^{2}\). It acts only inside the fiber of a fixed oriented absolute-amplitude block. This gives item (1).

Likewise, the present deeper sector introduces no new angular variable, no anisotropic locking term, and no quotient on constant values of \(\phi_{0}\). The minimizing branch of Theorem \ref{thm:branch} still carries arbitrary constant \(\phi_{0}\in\mathbb R/2\pi\mathbb Z\). Therefore the global \(O(2)\) vacuum angle remains residual, giving item (2).
\end{proof}

\begin{corollary}[When a later verdict revision would be honest]
\label{cor:revision}
The current scale-orbit and residual-angle verdicts may be revised only if some later still-deeper sector supplies a genuine quotient identifying distinct positive gap magnitudes and hence distinct positive moduli, or a genuine vacuum-angle locking or angle quotient beyond the internal pair-fiber structure of the present note.
\end{corollary}

\begin{proof}
This is the contrapositive reading of Theorem \ref{thm:verdict}. Without genuinely new quotient or locking data, the present verdicts remain.
\end{proof}

\begin{remark}
Corollary \ref{cor:revision} is intentionally strict. The present note does make the polarity labels more concrete by deriving them from the ordering of unequal positive pairs, but it still does not identify distinct positive moduli and it still does not lock the global \(O(2)\) vacuum angle.
\end{remark}

\section{Claim Boundary}

This note does not claim that the unequal positive-pair sector \(\PairSector\) is the unique final substrate completion below the oriented absolute-amplitude layer. It does not claim that the midpoint or offset data inside a pair-fiber class are already physically redundant, and it does not claim that every conceivable still-deeper completion must factor through the particular unequal positive-pair realization written here. It also does not claim that the positive scale orbit or the global \(O(2)\) vacuum angle have already been removed by a deeper redundancy or locking mechanism.

Its claim is narrower. The positive absolute-amplitude block and the discrete polarity data feeding the amplitude-sector origin note are no longer primitive either. They are generated by the explicit gap/order map \(\PairMap\) from a still-deeper sector of unequal positive pairs. Through the already-frozen origin map, amplitude lift, and setup map, that deeper sector reproduces the same signed amplitudes, the same readout block, the same ordered kinetic pair, the same radial anchor pair, the same positive orbit modulus, the same raw-orbit lift, and the same downstream package. The only quotient visible at the present stage is the internal pair-fiber quotient collapsing different deeper pair realizations of the same oriented absolute-amplitude data, and that quotient does not change the current scale-orbit or residual-angle verdicts.

\section{Conclusion}

The next honest manuscript step below the amplitude-sector note was to stop leaving the positive absolute amplitudes and discrete polarity data as the deepest recorded layer. This note now does that in the least-drift way. The absolute amplitudes are generated as the magnitudes of six unequal positive-pair gaps, and the polarity data are generated as the order orientation of those same gaps.

The positive result is exact and limited. Every oriented absolute-amplitude datum admits explicit lifts to the unequal positive-pair sector, parameterized by positive midpoint data. Composed with the frozen origin map, the present deeper sector reproduces the same signed amplitudes as the six pair gaps themselves. Composed further with the frozen amplitude lift and frozen setup map, it reproduces the same primitive class, the same readout block, the same ordered kinetic pair, the same radial anchor pair, the same positive orbit modulus, the same raw-orbit lift, the same density/current block, the same primitive bridge-order block, the same phase-neutral minimizing branch, the same canonical profile, and the same dressed bridge map.

The scope is equally clear. The only new quotient visible at this level is the internal pair-fiber quotient identifying different unequal positive-pair realizations of the same oriented absolute-amplitude datum. It does not identify distinct positive gap magnitudes, it does not identify distinct positive moduli, and it does not lock or quotient the global \(O(2)\) vacuum angle. Therefore the disciplined current verdict remains in force: the positive scale orbit is still a canonical-normalization orbit rather than a proved deeper redundancy, and the vacuum angle remains residual. The next honest open burden, if the derivational line continues, is therefore deeper again: a still more primitive origin, anchoring, or selection principle for the midpoint/fiber data and perhaps for the unequal positive-pair sector itself, with the verdicts revisited there only if genuinely new positive-modulus quotient or angle-locking structure appears.

\end{document}
